% from ega4-2-fr.tex
\oldpage[IV-2]{43}
$\mathcal N^{(x)}=\operatorname{Ker}(v)$, on a
$\mathcal N^{(x)}|U=\widetilde N$, donc
$x\notin\operatorname{Ass}(\mathcal N^{(x)})$. Il en résulte que
l'homomorphisme
$\mathcal F\to\bigoplus_{x\in\operatorname{Ass}(\mathcal F)}
\mathcal F^{(x)}$ est injectif, car son noyau $\mathcal H$ est contenu dans
tous les $\mathcal N^{(x)}$, donc $\operatorname{Ass}(\mathcal H)$ est contenu
dans l'intersection des $\operatorname{Ass}(\mathcal N^{(x)})$, qui est vide~;
par suite (3.1.5), $\mathcal H=0$. Compte tenu de ce que
$\operatorname{Ass}(\mathcal F)$ est localement fini (3.1.6), il est clair que
$(\mathcal F^{(x)})_{x\in\operatorname{Ass}(\mathcal F)}$ est une
décomposition irredondante réduite de $\mathcal F$ vérifiant les conditions de
l'énoncé. La caractérisation de $\mathcal F^{(x)}$ lorsque
$\overline{\{x\}}$ n'est pas immergé résulte de Bourbaki,
\emph{Alg. comm.}, chap.~IV, \S~2, n\textsuperscript{o}~3, prop.~5, la question
étant locale, et compte tenu de (I, 1.6.7).

\begin{corollary}[3.2.7]
\label{IV.2.3.2.7-fr}
Sous les hypothèses de (3.2.6), si $\mathcal F$ n'a pas de cycle premier
associé immergé, il n'existe qu'une seule décomposition irredondante réduite de
$\mathcal F$.
\end{corollary}

\begin{corollary}[3.2.8]
\label{IV.2.3.2.8-fr}
Soient $X$ un préschéma noethérien, $\mathcal F$ un $\mathcal O_X$-Module
cohérent. Il existe une filtration finie $(\mathcal F_i)_{0\le i\le n}$ de
$\mathcal F$ telle que $\mathcal F_0=\mathcal F$, $\mathcal F_n=0$, formée de
$\mathcal O_X$-Modules cohérents et tels que les quotients
$\mathcal F_i/\mathcal F_{i+1}$ soient nuls ou irredondants, et
$\operatorname{Ass}(\mathcal F_i/\mathcal F_{i+1})
\subset\operatorname{Ass}(\mathcal F)$.
\end{corollary}

En effet, $\mathcal F$ est isomorphe à un sous-$\mathcal O_X$-Module d'une
somme directe finie $\bigoplus_{j=1}^n\mathcal G_j$, où les $\mathcal G_j$
sont irredondants et cohérents (3.2.6)~; comme tout sous-$\mathcal O_X$-Module
quasi-cohérent de $\mathcal G_j$ est nul ou irredondant (3.1.7), les
$\mathcal F_i=\mathcal F\cap(\bigoplus_{j=1}^{n-i}\mathcal G_j)$ répondent à
la question, $\mathcal F_i/\mathcal F_{i+1}$ étant isomorphe à un
sous-$\mathcal O_X$-Module cohérent de $\mathcal G_{n-i}$.

\subsection{Relations avec la platitude}

\begin{proposition}[3.3.1]
\EGAReaderTarget{ega4.c04.prop.3-3-1}{IV.2.3.3.1-fr}
Soient $f:X\to Y$ un morphisme, $\mathcal F$ un $\mathcal O_X$-Module
quasi-cohérent et $f$-plat, $\mathcal E$ un $\mathcal O_Y$-Module
quasi-cohérent. Si, pour tout $y\in Y$, on pose
$\mathcal F_y=\mathcal F\otimes_{\mathcal O_Y}k(y)$, on a
\end{proposition}

\begin{env}[3.3.1.1]
\label{IV.2.3.3.1.1-fr}
\[
\operatorname{Ass}(\mathcal E\otimes_Y\mathcal F)
\supset\bigcup_{y\in\operatorname{Ass}(\mathcal E)}
\operatorname{Ass}(\mathcal F_y),
\]
\end{env}

et les deux membres sont égaux si $Y$ est localement noethérien.

(Bien entendu, $\mathcal F_y$ est un faisceau sur la fibre $f^{-1}(y)$, et on
identifie cette fibre à un sous-espace de $X$ (I, 3.6.1).) La question étant
locale sur $X$ et sur $Y$, on est ramené au cas où $X$ et $Y$ sont affines,
et la proposition est alors démontrée dans Bourbaki, \emph{Alg. comm.},
chap.~IV, \S~2, n\textsuperscript{o}~6, th.~2.

\begin{corollary}[3.3.2]
\label{IV.2.3.3.2-fr}
Soient $Y$ un préschéma localement noethérien sans cycles premiers associés
immergés, $f:X\to Y$ un morphisme, $\mathcal F$ un $\mathcal O_X$-Module
quasi-cohérent et $f$-plat. Alors, pour tout
$x\in\operatorname{Ass}(\mathcal F)$, $f(x)$ est point maximal de $Y$.
\end{corollary}

Il suffit d'appliquer (3.3.1) avec $\mathcal E=\mathcal O_Y$, puisque
$\operatorname{Ass}(\mathcal O_Y)$ est par hypothèse l'ensemble des points
maximaux de $Y$.

\begin{corollary}[3.3.3]
\label{IV.2.3.3.3-fr}
Sous les hypothèses de (3.3.1), supposons en outre que $X$ et $Y$ soient
localement noethériens, $\mathcal E$ et $\mathcal F$ cohérents. Alors les
conditions suivantes sont équivalentes~:
\begin{enumerate}
\item[\emph{a)}] $\mathcal E\otimes_Y\mathcal F$ est sans cycle premier
associé immergé.
\item[\emph{b)}] Pour tout point
$y\in\operatorname{Ass}(\mathcal E)\cap f(\operatorname{Supp}(\mathcal F))$,
$\overline{\{y\}}$ est un cycle premier associé non immergé de $\mathcal E$ et
$\mathcal F_y$ est sans cycle premier associé immergé.
\end{enumerate}
\end{corollary}

